% RESEARCH_DRAFT_SHA256: d17f867415b96471128139a6ee35e1bbe4c320f954e2cc3ab7eabfaaaac9be2e
% Layout and source-reference wrappers expanded for the web renderer.
\documentclass[11pt]{article}
\usepackage[margin=25mm]{geometry}
\usepackage{amsthm}
\usepackage{fontspec,unicode-math}
\setmainfont{Latin Modern Roman}
\setmathfont{Latin Modern Math}
\usepackage{amsmath,amssymb,mathtools}
\usepackage{xurl,hyperref}
\hypersetup{colorlinks=true,urlcolor=blue,linkcolor=blue}
\setcounter{secnumdepth}{0}
\setlength{\parindent}{0pt}
\setlength{\parskip}{5pt}
\setlength{\emergencystretch}{3em}
\newcommand{\sref}[2]{\hyperlink{source:#1:#2}{[S#2]}}
\newcommand{\cataloguescope}[1]{\paragraph{Source question.}#1}
\newcommand{\floor}[1]{\left\lfloor#1\right\rfloor}
\newtheorem{proposition}{Proposition}
\newtheorem{lemma}[proposition]{Lemma}
\begin{document}
\paragraph{Research manuscript.} OpenAI \texttt{gpt-6-astra}, reasoning \texttt{ultra}; 11 October 2026.
\paragraph{Preserved source body.} The complete catalogue mathematical text and sources follow with layout and source links expanded, including source qualifications and its historical preparation notes. New research begins at the section headed Research attempt.

% UnsolvedMath ID 30000608; source catalogue code OWR-1386-015
\section{Complexity of Dense Perfect Hypergraph Matching}\label{30000608}
{\small UnsolvedMath ID 30000608 · Computer Science}
\subsection{Definitions and mathematical statement}

\paragraph{Shared notation.}Standard number systems are denoted by $\mathbb N,\mathbb Z,\mathbb Q,\mathbb R,\mathbb C$; index ranges, computational models and specialized notation follow the cited problem source and the definitions below.


A $k$-uniform hypergraph has edges that are $k$-element vertex subsets. A matching is a collection of pairwise disjoint edges. Its minimum $(k-1)$-codegree is the minimum number of edges containing a given $(k-1)$-subset.

A graph matching is a set of edges with pairwise disjoint endpoints, and it is perfect when every vertex is incident to one chosen edge. A weighted matching objective sums the chosen edge weights, unless an ordered objective is specified.
\paragraph{Statement recorded in the supplied source.}
For integers $k\geq2$ and $1/k\leq c<1/2$, determine the complexity of deciding whether a $k$-uniform hypergraph $H$ with $|V(H)|\equiv0\pmod k$ and $\delta_{k-1}(H)\geq cn$ has a matching of size $n/k$. \href{https://doi.org/10.4171/owr/2006/48}{[S2]}
\paragraph{Source status.}
The supplied catalogue labels this record open. Its source-status metadata was checked against the supplied snapshot on 2026-10-11. This import does not independently establish current open status or mathematical correctness. The catalogue literature assessment is dated 2026-08-21.
\subsection{Short English statement}
Classify the difficulty of finding perfect matchings in uniform hypergraphs above the stated codegree threshold.
\subsection{Sources}
\begin{itemize}
\item[S1] Ulam AI and the UnsolvedMath Contributors, UnsolvedMath: A Curated Collection of Open Mathematics Problems, record 30000608 (OWR-1386-015). Catalogue curation and metadata: CC BY 4.0. The supplied source question is retained; definitions and the English summary are editorial additions.
\url{https://huggingface.co/datasets/ulamai/UnsolvedMath}.
\item[S2] Combinatorics, Probability and Computing (2007), Oberwolfach Reports, DOI 10.4171/owr/2006/48, pp. 2875--2952
\url{https://doi.org/10.4171/owr/2006/48}.
\textit{Underlying Oberwolfach Reports material retains its source-specific licence, including CC BY-SA 4.0 where applicable; see THIRD\_PARTY\_NOTICES.md.}
\end{itemize}
\paragraph{Preparation and review.}
Original question attributed to the cited source and UnsolvedMath contributors. Codex (GPT-6; exact serving revision not exposed) assisted with definitions, English summaries, source inspection and formatting. No new human mathematical authorship or independent proof verification is asserted. Canonical ID 30000608 and source code OWR-1386-015 checked against the complete supplied registry and existing repository statements; known flagged repetitions and obvious identical conjectures excluded. Full mathematical-equivalence/literature review has not been performed. Statement preparation and final review followed contribution prompts C and D; automated checks verify structure and identity, not open status.
\label{end:30000608}

\section{Research attempt}
\subsection{Literature, source context and scope}
The full original catalogue body and source qualifications are preserved above. The relevant primary workshop passages were read in the cached publisher report. Historical claims and subsequent attributed results are distinguished from the mathematical work below. A restricted example, standard lemma or known theorem is not counted as new progress on the entire original target. Missing assumptions and repeated or rhetorical extracts are not silently turned into different solved problems.

\subsection{A dense no-instance and the distinction between existence and decision}
The report asks for the complexity of perfect matching under minimum codegree $cn$, for fixed $k\ge3$ and $1/k\le c<1/2$. The later Han--Keevash primary paper explicitly treats fixed $k$ and states a polynomial-time constructive algorithm even at codegree at least $n/k-O_k(\log n)$; its introduction also attributes the earlier decision classification to Han. Thus the historical extract does not incorporate that literature. We do not claim its algorithm as our result. Instead we test whether density alone could make every allowed instance a yes-instance.
\begin{proposition}
Let $k\mid n$, partition the vertices as $A\cup B$ with $|A|$ odd, and take exactly the $k$-sets meeting $A$ in an even number of vertices. This hypergraph has no perfect matching and has minimum codegree at least $\min(|A|,|B|)-(k-1)$.
\end{proposition}
\begin{proof}
Every edge covers an even number of vertices in $A$, so the union of any disjoint collection of edges also covers an even number there. It cannot cover all of the odd set $A$, proving absence of a perfect matching. For a fixed $(k-1)$-set $S$, if $|S\cap A|$ is odd, all extensions in $A\setminus S$ and only those have even intersection with $A$. If it is even, all extensions in $B\setminus S$ and only those qualify. In either case at least $\min(|A|,|B|)-(k-1)$ vertices extend $S$ to an edge.
\end{proof}
Choose $|A|$ odd within one of $n/2$. Then the lower bound is at least $n/2-k$, which exceeds $cn$ for every fixed $c<1/2$ once $n$ is sufficiently large. Consequently the original density promise does not itself force existence. The source already discusses parity constructions of this kind; this proof is an explicit analysis of the failed density-only route, not a new obstruction theorem.
\subsection{Verification and remaining gaps}
For $k=3,n=12,|A|=5$, exhaustive edge and vertex-mask matching calculations give minimum codegree 4 and maximum matching size 3, although $n/k=4$. Thus a no-instance occurs even at the lower endpoint $c=1/3$. Deciding no and providing a certificate can still be polynomial-time, as the attributed literature states; hard instances and no-instances are not synonymous. No independently derived algorithm for the full promise class, complexity lower bound, or novelty claim is supplied. The external theorem was checked at the statement and scope level, not by verifying its entire proof.

\section{References and source attribution}
Primary workshop report: \url{https://doi.org/10.4171/owr/2006/48}. Relevant passages read: 2006/48 pp. 2929--2930, parity barriers and the promised perfect-matching complexity question.

\url{https://arxiv.org/pdf/2210.12643}. Read primary introduction pp. 1--3, fixed-k convention, Theorem 1.2 and its scope; it supplies a polynomial constructive algorithm at codegree n/k minus O_k(log n) and credits earlier decision results. Full proof not independently checked.

All retrievals were checked on 11 October 2026. Original question attribution remains in the preserved Sources section. External theorems are attributed and not claimed as new discoveries or fully verified merely by citation.


\paragraph{Integration review.} Codex (GPT-6; exact serving revision not exposed) reviewed the written arguments and expanded layout and source-reference wrappers. Repository metadata and the production web renderer were checked using Python 3.13 and JavaScript with MathJax 3.2.2. This model-assisted review is not independent human proof verification; no PDF compilation or formal proof checking is asserted.

\section{Completion Estimate}
\noindent\textbf{COMPLETION: 0\%}
No substantiated new progress toward the entire original target is claimed. The worked results and failed route are delimited above; a known result, illustrative model, or auxiliary obstruction is not counted as original-target completion.

\section{Final status and attribution}
\textbf{Attempt status: unresolved.} This describes this research attempt, not a new ruling on the source catalogue's status or on external literature claims.
The mathematics and this self-audit were generated by OpenAI \texttt{gpt-6-astra}, reasoning \texttt{ultra}, on 11 October 2026. Codex (GPT-6; exact serving revision not exposed) orchestrated the batch and contributed formatting/repository review. No human mathematical authorship, independent human proof checking, or formal verification is claimed. The historical preparation paragraphs in the reproduced catalogue body concern that body only.

Contribution prompt A was applied to the complete definition and supplied source record, followed by the adversarial checks of prompt B. Local reproduction files are \texttt{batch016/checks.py}, \texttt{computations.json}, and the per-ID audit JSON. They were not public links at preparation time. Numerical checks are finite and do not prove unrestricted results. No LaTeX engine was found locally; compilation was not performed by this research agent.

\end{document}
